% exact vector join: batches move across disk, host, and GPU; one task joins one batch pair
\begin{figure*}
    \centering
    \definecolor{vjR}{HTML}{3B6FB6}
    \definecolor{vjS}{HTML}{E07B39}
    \definecolor{vjD}{HTML}{6A4C93}
    % one batch: light box, darker stripes for its rows (vectors)
    \newcommand{\vjbatch}[3]{% name, position, color
        \node[draw=#3!80!black, fill=#3!12, rounded corners=1.5pt, minimum width=0.8cm,
            minimum height=0.5cm, inner sep=0pt] (#1) at (#2) {};
        \foreach \dy in {-0.12,0,0.12}
        \draw[#3!75, line width=1.3pt, line cap=round]
        ([xshift=0.12cm, yshift=\dy cm]#1.west) -- ([xshift=-0.12cm, yshift=\dy cm]#1.east);
    }
    % a batch in a memory tier, labeled inside
    \newcommand{\vjslot}[4]{% name, position, fill, label
        \node[draw=black!45, fill=#3, rounded corners=1.5pt, minimum width=0.55cm,
            minimum height=0.38cm, inner sep=0pt, font=\tiny] (#1) at (#2) {#4};
    }
    % a table on disk: header band and two row lines
    \newcommand{\vjtable}[2]{% position, color
        \begin{scope}[shift={(#1)}]
            \filldraw[fill=white, draw=#2!80!black] (-0.3,-0.22) rectangle (0.3,0.22);
            \fill[#2!60] (-0.3,0.1) rectangle (0.3,0.22);
            \draw[#2!80!black] (-0.3,0.1) -- (0.3,0.1);
            \draw[#2!40] (-0.3,-0.02) -- (0.3,-0.02) (-0.3,-0.12) -- (0.3,-0.12);
        \end{scope}
    }
    % 3x5 distance tile with the cells that survive the local reduction filled in
    \newcommand{\vjmini}[4]{% name, x, y, kept cells as row/col
        \begin{scope}[shift={(#2,#3)}]
        \foreach \ri in {0,1,2} {\foreach \ci in {0,...,4} {
            \filldraw[fill=vjD!6, draw=black!30, line width=0.3pt]
            ({-0.4+\ci*0.16},{0.08-\ri*0.16}) rectangle ++(0.16,0.16);}}
        \foreach \ri/\ci in {#4} {
            \filldraw[fill=vjD!80, draw=black!30, line width=0.3pt]
            ({-0.4+\ci*0.16},{0.08-\ri*0.16}) rectangle ++(0.16,0.16);}
        \end{scope}
        \node[minimum width=0.8cm, minimum height=0.48cm, inner sep=0pt] (#1) at (#2,#3) {};
    }
    \begin{tikzpicture}[
        font=\scriptsize\sffamily,
        panel/.style={fill=black!5, draw=black!20, rounded corners=4pt},
        ptitle/.style={anchor=west, font=\small\sffamily\bfseries},
        head/.style={align=center},
        flow/.style={-{Latex[length=1.6mm]}, black!70, line width=0.6pt},
        tier/.style={-{Latex[length=1.4mm]}, black!60, line width=0.7pt},
        sel/.style={draw=black!60, fill=white, rounded corners=2pt, minimum width=1.55cm,
        minimum height=0.75cm, align=center},
        op/.style={draw=black!60, fill=white, rounded corners=2pt, minimum width=2.5cm,
        minimum height=0.8cm, align=center},
        sub/.style={font=\tiny\sffamily},
    ]

        % ---------- zoom from GPU memory into the GPU work: dashed lines between panel corners ----------
        \draw[black!45, dashed] (4.5,2.15) -- (5.3,5.85) (4.5,0.15) -- (5.3,0.15);

        % ---------- memory tiers: disk on top, host in the middle, GPU at the bottom ----------
        \draw[panel] (0,4.6) rectangle (4.5,5.85);
        \node[ptitle] at (0.1,5.6) {\faHdd~Disk};
        \draw[panel] (0,2.75) rectangle (4.5,4.0);
        \node[ptitle] at (0.1,3.75) {\faMemory~Host memory};
        \draw[panel] (0,0.15) rectangle (4.5,2.15);
        \node[ptitle] at (0.1,1.9) {\faMicrochip~GPU memory};

        % disk: the two input tables
        \vjtable{0.6,5.05}{vjR}
        \node[anchor=west] at (0.95,5.05) {$R$ (probe)};
        \vjtable{2.6,5.05}{vjS}
        \node[anchor=west] at (2.95,5.05) {$S$ (corpus)};

        % host: pinned or spilled batches and intermediate results
        \vjslot{h1}{0.5,3.2}{vjR!15}{}
        \vjslot{h2}{1.1,3.2}{vjS!15}{}
        \vjslot{h3}{1.7,3.2}{vjD!15}{}
        \node[sub, anchor=west, align=left] at (2.1,3.2) {pinned or spilled batches,\\intermediate results};

        % GPU: both relations split into batches; the pair of the zoomed task is highlighted
        \foreach \x/\rl/\sl in {0.55/1/1, 1.25/2/2, 3.65/p/q} {
            \vjslot{r\rl}{\x,1.25}{vjR!15}{$R_\rl$}
            \vjslot{s\sl}{\x,0.6}{vjS!15}{$S_\sl$}}
        \vjslot{ri0}{2.55,1.25}{vjR!45}{$R_i$}
        \vjslot{sj0}{2.55,0.6}{vjS!45}{$S_j$}
        \foreach \x in {1.9,3.1} {
            \node at (\x,1.25) {$\cdots$};
            \node at (\x,0.6) {$\cdots$};}

        % moves between tiers: load goes down, cuCascade spills back up
        \foreach \yb/\yt in {4.0/4.6, 2.15/2.75} {
            \draw[tier] (1.9,\yt) -- (1.9,\yb) node[midway, left, sub] {load};
            \draw[tier] (2.5,\yb) -- (2.5,\yt) node[midway, right, sub] {spill (cuCascade)};}

        % ---------- batch pairs on the GPU ----------
        \draw[panel] (5.3,0.15) rectangle (17.55,5.85);
        \node[ptitle] at (5.4,5.6) {\faMicrochip~GPU: up to $T$ batch pairs in parallel};

        % vertical lines between steps
        \foreach \x in {6.75,9.25,13.05}
        \draw[black!25] (\x,0.35) -- (\x,5.25);

        % step headers: titles on one baseline, subtitles on another
        \foreach \hx/\htitle/\hsub in {
            6.02/{\textbf{Input}}/{one pair},
            8.0/{\textbf{Step 1}: Distance comp.}/{matrix multiplication},
            11.15/{\textbf{Step 2a}: Local reduction}/{within a pair},
            15.3/{\textbf{Step 2b}: Global reduction}/{across pairs}} {
            \node[anchor=base] at (\hx,4.95) {\htitle};
            \node[anchor=base, sub] at (\hx,4.7) {\hsub};}

        % ---------- input pair ----------
        % faint copies behind: other tasks work on their own pairs at the same time
        \foreach \sh in {0.16,0.08} {
            \node[draw=vjS!35, fill=vjS!6, rounded corners=1.5pt, minimum width=0.8cm,
                minimum height=0.5cm, inner sep=0pt] at ({6.15+\sh},{3.35+\sh}) {};
            \node[draw=vjR!35, fill=vjR!6, rounded corners=1.5pt, minimum width=0.8cm,
                minimum height=0.5cm, inner sep=0pt] at ({6.15+\sh},{2.25+\sh}) {};}
        \vjbatch{sj}{6.15,3.35}{vjS}
        \vjbatch{ri}{6.15,2.25}{vjR}
        \node[left=1pt of sj] {$S_j$};
        \node[left=1pt of ri] {$R_i$};

        % ---------- step 1: R_i times S_j^T gives the distances ----------
        % rows of R_i on the left, columns of S_j^T on top, one row and column highlighted
        \foreach \ri in {0,...,3} {
            \pgfmathsetmacro{\c}{\ri==1 ? 85 : 35}
            \filldraw[fill=vjR!\c, draw=white, line width=0.3pt]
            (7.15,{2.8-\ri*0.2}) rectangle ++(0.2,0.2);}
        \foreach \ci in {0,...,5} {
            \pgfmathsetmacro{\c}{\ci==2 ? 85 : 35}
            \filldraw[fill=vjS!\c, draw=white, line width=0.3pt]
            ({7.45+\ci*0.2},3.1) rectangle ++(0.2,0.2);}
        \foreach \ri in {0,...,3} {\foreach \ci in {0,...,5} {
            \pgfmathtruncatemacro{\s}{mod(3*\ri+7*\ci+1,5)*15+10}
            \filldraw[fill=vjD!\s, draw=white, line width=0.3pt]
            ({7.45+\ci*0.2},{2.8-\ri*0.2}) rectangle ++(0.2,0.2);}}
        \draw[black!50] (7.45,2.2) rectangle (8.65,3.0);
        \draw[black, line width=0.7pt] (7.85,2.6) rectangle (8.05,2.8);
        \node[sub] at (7.25,2.02) {$R_i$};
        \node[sub] at (8.05,3.45) {$S_j^\top$};
        \node[sub] at (8.05,2.02) {distances};
        \node[minimum width=1.2cm, minimum height=0.8cm, inner sep=0pt] (tile) at (8.05,2.6) {};
        \draw[flow] (sj.east) -- (7.4,3.2);
        \draw[flow] (ri.east) -- (7.1,2.5);

        % ---------- step 2a: selection ----------
        \node[sel] (topk) at (10.15,3.85) {top-$k$\\select};
        \node[sel] (eps)  at (10.15,1.35) {$\varepsilon$\\filter};
        \draw[flow] ([xshift=0.22cm]tile.east) -- (topk.west);
        \draw[flow] ([xshift=0.22cm]tile.east) -- (eps.west);

        % local results of one pair (dark cells survive), k = 2; global and threshold cells agree with per-row
        \vjmini{prmini}{11.95}{3.85}{0/1,0/3,1/0,1/2,2/2,2/4}
        \vjmini{glmini}{11.95}{2.6}{0/1,1/0}
        \vjmini{thmini}{11.95}{1.35}{0/1,1/0,1/2,1/4}
        \node[sub, inner sep=0pt, above=0.5pt of prmini] {\textbf{Per-row}: $k$ per row};
        \node[sub, inner sep=0pt, below=1.5pt of glmini] {\textbf{Global}: $k$ per pair};
        \node[sub, inner sep=0pt, above=0.5pt of thmini] {\textbf{Threshold}: $\le \varepsilon$};
        \draw[flow] (topk.east) -- (prmini.west);
        % global keeps the k best of the per-row result
        \draw[flow] (prmini.south) -- (glmini.north) node[midway, left, sub] {keep $k$ best};
        \draw[flow] (eps.east) -- (thmini.west);

        % ---------- step 2b: global reduction ----------
        \node[op] (prop) at (14.45,3.85)
            {\texttt{VECTOR\_TOPK\_MERGE}\\{\tiny merge $q$ partials per $R_i$}};
        \node[op] (glop) at (14.45,2.6)
            {\texttt{TOP\_N} $\rightarrow$ \texttt{MERGE\_TOP\_N}\\{\tiny top-$k$ of all candidates}};
        \draw[flow] (prmini.east) -- (prop.west);
        \draw[flow] (glmini.east) -- (glop.west);

        % ---------- outputs ----------
        \node[align=center] (prout) at (16.75,3.85) {$k$ per\\probe row};
        \node[align=center] (glout) at (16.75,2.6) {$k$ pairs};
        \node[align=center] (thout) at (16.75,1.35) {all pairs\\within $\varepsilon$};
        \draw[flow] (prop.east) -- (prout.west);
        \draw[flow] (glop.east) -- (glout.west);
        % threshold needs no global step: the union of local results is the answer
        \draw[flow] (thmini.east) -- (thout.west);
    \end{tikzpicture}
    \captionsetup{textfont=normalfont}
    \caption{
        Exact vector join. The GPU joins both tables in batches, which move between disk, host, and GPU memory.
    }
    \label{fig:exact-join}
\end{figure*}
